% Generated from the audited Markdown source; responsible author: Carptopus.
\documentclass[11pt]{article}
\usepackage[a4paper,margin=28mm]{geometry}
\usepackage[T1]{fontenc}
\usepackage[utf8]{inputenc}
\usepackage{lmodern}
\usepackage{microtype}
\usepackage{amsmath,amssymb,amsthm,mathtools}
\usepackage{xcolor}
\usepackage[colorlinks=true,linkcolor=blue!55!black,citecolor=blue!55!black,urlcolor=blue!65!black]{hyperref}
\usepackage{fancyhdr}

\newtheorem{theorem}{Theorem}[section]
\newtheorem{lemma}[theorem]{Lemma}
\newtheorem{corollary}[theorem]{Corollary}
\newtheorem{proposition}[theorem]{Proposition}
\numberwithin{equation}{section}

\pagestyle{fancy}
\fancyhf{}
\fancyhead[L]{Carptopus}
\fancyhead[R]{Measurable lattice-path matroids}
\fancyfoot[C]{\thepage}
\setlength{\headheight}{14pt}
\setlength{\parindent}{0pt}
\setlength{\parskip}{0.45em}
\setlength{\emergencystretch}{4em}

\title{Minorizing measures and truncations\\
of measurable lattice-path matroids}
\author{Carptopus\\\href{mailto:carptopus@163.com}{carptopus@163.com}}
\date{24 August 2026}

\hypersetup{
  pdftitle={Minorizing measures and truncations of measurable lattice-path matroids},
  pdfauthor={Carptopus},
  pdfsubject={Exact minorizing-measure relaxations for measurable lattice-path matroids},
  pdfkeywords={measurable matroid, lattice-path matroid, minorizing measure, weak-star topology, truncation, extreme point, continuous linear programming}
}

\begin{document}
\maketitle

\begin{center}
\fcolorbox{black!35}{black!4}{\parbox{0.88\textwidth}{\centering\small
Public Beta v0.1-beta. The mathematical claims have passed a project-internal
zero-trust manuscript audit; external review is pending.}}
\end{center}

\begin{abstract}


We study the minorizing-measure sets of the measurable lattice-path matroids recently
introduced by Bérczi, Imolay, and Schweitzer. Such a matroid is specified by a closed
set of times and by lower and upper cumulative paths. We first prove that its set of
basic minorizing measures is exactly the weak-star compact corridor of fractional path
densities. The proof is constructive: after moving a fractional path slightly toward
a genuine basis path, we purify it on the components of the strict interior by a
countable endpoint-preserving bang--bang construction.

We then determine the full minorizing-measure set. A density belongs to it if and only
if all interval increments satisfy the inequalities forced by the two boundary paths.
The nontrivial direction is proved by an explicit residual-capacity time change that
extends the density to a basic minorizing measure. As a consequence, the basic
minorizing measures of every real truncation are described by the same interval
inequalities together with one total-mass equation. We also record the contact-set
criterion for extreme points, show directly that every extreme point is weak-star
exposed, and give a bounded laminar-composition application.

The finite-dimensional shadows of the corridor and interval descriptions are classical
lattice-path and positroid polytope phenomena. The contribution here is their exact
atomless measurable counterpart, including arbitrary closed time sets and real
truncation levels.

\end{abstract}

\section{Introduction}


Measurable matroids provide a measure-theoretic analogue of finite matroids on standard
atomless measure spaces. Bérczi, Imolay, and Schweitzer established equivalent
independent-set, basis, rank, and closure descriptions, introduced measurable versions
of standard matroid classes and operations, and proved that independent-set and basis
measures are weak-star dense in the corresponding minorizing-measure sets \cite{berczi2026}. Their
examples include measurable lattice-path matroids: a basis is represented by a
zero--one derivative whose cumulative path stays between two prescribed boundaries.

The purpose of this paper is to make the associated fractional geometry explicit.
There are two distinct convex sets:

\begin{enumerate}
\item the basic minorizing measures, which have the full matroid rank as total mass; and
\item all minorizing measures, whose total mass is allowed to vary.

\end{enumerate}

For the first set, the expected candidate is the obvious path corridor. In a finite
lattice-path matroid the analogous base-polytope description is classical \cite{bidkhori2012,ohxiang2022}. In the
atomless setting, however, it is not enough to observe that every basis lies in the
corridor: one must prove that every fractional corridor path is a weak-star limit of
genuine basis paths. We give a self-contained purification proof that remains valid for
an arbitrary closed set of constrained times, including infinitely many complementary
gaps and Cantor-type contact sets.

The full minorizing-measure set is subtler. Prefix lower bounds no longer make sense
without fixed total mass. They are replaced by interval inequalities

\nopagebreak[4]
\[
F(t)-F(s)\le b(t)-a(s)\qquad(s\le t),
\]

where $F$ is the cumulative density. We prove that these inequalities are not merely
necessary: they are sufficient. The proof builds, from the unused pointwise capacity
$1-f$, a complementary path that raises the density to a point in the basic
minorizing-measure set. This yields the following informal form of the main result:

\begin{quote}\itshape
The complete independence relaxation of a measurable lattice-path matroid is its
interval relaxation, and every real truncation is obtained by taking a mass slice.
\end{quote}


The extreme-point criterion for the fixed-mass corridor is closely related to classical
extreme-point results in separated continuous linear programming \cite{anderson1983,andersonnash1987}. We include a
short direct proof because it identifies the exact contact set and gives a concrete
exposing functional without relying on a claimed theorem-level equivalence. We do not
claim a new general bang--bang or continuous-control theorem: related purification
mechanisms have a substantial history \cite{amarcellina1994,carlota2016}.

The paper is organized as follows. Section 2 gives the setup. Section 3 proves the
basic-minorizing-measure identity. Section 4 determines the full minorizing-measure set
and all truncation slices. Section 5 treats extreme and exposed points. Section 6 gives
a finite-depth laminar-composition application and carefully separates it from the main
contribution. Section 7 records the scope and remaining open boundary.

\section{Setup}


\subsection{Minorizing measures}


Let $(J,\Sigma,\mu,r)$ be a measurable matroid in the sense of \cite{berczi2026}. Throughout,
$\alpha$ denotes a finite nonnegative charge on $(J,\Sigma)$. Its set of minorizing
measures and its set of basic minorizing measures are

\nopagebreak[4]
\[
\mathsf{mm}(r)
=\{\alpha:0\le \alpha(X)\le r(X)\text{ for every }X\in\Sigma\},
\]

and

\nopagebreak[4]
\[
\mathsf{bmm}(r)
=\{\alpha\in\mathsf{mm}(r):\alpha(J)=r(J)\}.
\]

Since $r\le\mu$, every such measure is absolutely continuous with respect to $\mu$ and
has a density in the order interval $0\le f\le1$. We use the weak-star topology of
$L^\infty(\mu)=(L^1(\mu))^*$. Proposition 5.6 and Corollary 5.7 of \cite{berczi2026} state that
independent-set measures are weak-star dense in $\mathsf{mm}(r)$ and basis measures are
weak-star dense in $\mathsf{bmm}(r)$.

\subsection{Measurable lattice-path data}


Let $K\subseteq[0,1]$ be closed and contain $0$ and $1$. Fix $0\le R\le1$ and let

\nopagebreak[4]
\[
a,b:K\longrightarrow[0,R]
\]

be nondecreasing $1$-Lipschitz functions satisfying

\nopagebreak[4]
\[
a(0)=b(0)=0,\qquad a(1)=b(1)=R,\qquad a(t)\le b(t).
\]

Assume that the basis family

\nopagebreak[4]
\[
\mathcal B_K(a,b)
=\left\{B\subseteq[0,1]:
\lambda(B)=R,
\ a(t)\le\lambda(B\cap[0,t])\le b(t)\quad(t\in K)
\right\}
\]

is nonempty. By Theorem 3.13 of \cite{berczi2026}, it is the basis family of a measurable matroid,
which we denote by $M_K(a,b)$. Its rank function is denoted by $r$.

For $f\in L^\infty([0,1])$, set

\nopagebreak[4]
\[
F_f(t)=\int_0^t f(x)\,dx.
\]

We use two convex relaxations. The fixed-mass path corridor is

\nopagebreak[4]
\[
\mathcal P_K(a,b)
=\left\{f:
0\le f\le1,\ F_f(1)=R,
\ a(t)\le F_f(t)\le b(t)\quad(t\in K)
\right\}.
\]

The interval relaxation is

\nopagebreak[4]
\[
\mathcal I_K(a,b)
=\left\{f:
0\le f\le1,
\ F_f(t)-F_f(s)\le b(t)-a(s)
\quad(s,t\in K,\ s\le t)
\right\}.
\]

Both sets are weak-star compact. Indeed, the order interval $0\le f\le1$ is weak-star
compact, and every cumulative evaluation $f\mapsto F_f(t)$ is weak-star continuous.

\section{The basic minorizing-measure corridor}


The first main result identifies the fixed-mass relaxation exactly.

\begin{theorem}[corridor identity]
\label{thm:corridor}

For every nonempty measurable lattice-path matroid $M_K(a,b)$,

\nopagebreak[4]
\[
\mathsf{bmm}(M_K(a,b))
=\{f\lambda:f\in\mathcal P_K(a,b)\}.
\]

The inclusion from left to right follows immediately from the weak-star density of
basis measures. The reverse inclusion is the content of the following purification
argument.

\end{theorem}

\noindent\emph{Proof.}

\subsection{Completing the obstacles across the gaps}


Let $(s,t)$ be a component of $[0,1]\setminus K$. Define on this gap

\nopagebreak[4]
\[
\ell(x)=\max\{a(s),a(t)-(t-x)\},
\]

and

\nopagebreak[4]
\[
u(x)=\min\{b(s)+x-s,b(t)\}.
\]

On $K$, put $\ell=a$ and $u=b$. These definitions give continuous, nondecreasing,
$1$-Lipschitz functions on $[0,1]$. Moreover, a nondecreasing $1$-Lipschitz path $F$
satisfies the constraints at the endpoints of a gap if and only if

\nopagebreak[4]
\[
\ell(x)\le F(x)\le u(x)
\]

throughout the gap. Thus $\mathcal P_K(a,b)$ is exactly the set of derivatives of
absolutely continuous paths $F$ such that

\nopagebreak[4]
\[
F(0)=0,\qquad F(1)=R,\qquad
0\le F'\le1,\qquad \ell\le F\le u.
\]

Because $\mathcal B_K(a,b)$ is nonempty, there is a path $H_0$ in this corridor with
$H_0'\in\{0,1\}$ almost everywhere.

\subsection{A local endpoint-preserving purification}


We first record an elementary lemma.

\begin{lemma}
\label{lem:local-purification}

Let $Y$ be absolutely continuous on a compact interval $[p,q]$ and satisfy
$0\le Y'\le1$ almost everywhere. For every $\eta>0$, there exists an absolutely
continuous $W$ such that

\nopagebreak[4]
\[
W(p)=Y(p),\qquad W(q)=Y(q),\qquad W'\in\{0,1\}\text{ a.e.},
\]

and $\|W-Y\|_\infty<\eta$.

\end{lemma}

\begin{proof}

Choose a mesh of width smaller than $\eta$. On a mesh interval $[r,s]$, let
$\Delta=Y(s)-Y(r)$. Since $0\le\Delta\le s-r$, define $W$ to move with slope $1$
for time $\Delta$ and with slope $0$ for the rest of the interval. The new and old
paths agree at all mesh points and both remain between the same two endpoint values on
each mesh interval. The uniform error is therefore smaller than the mesh width.
\end{proof}


\subsection{Countable purification inside the corridor}


Fix $f\in\mathcal P_K(a,b)$ and write $F=F_f$. For $0<\delta<1$, set

\nopagebreak[4]
\[
F_\delta=(1-\delta)F+\delta H_0.
\]

Let

\nopagebreak[4]
\[
Z=\{x:F_\delta(x)=\ell(x)\text{ or }F_\delta(x)=u(x)\}.
\]

This is closed. If, for example, $F_\delta=\ell$ at a point, then both
$F-\ell$ and $H_0-\ell$ are nonnegative there and their strict convex combination is
zero. Hence $F=H_0=\ell$ at that point. The analogous statement holds on the upper
contact set. The zero-set derivative lemma for absolutely continuous functions now
implies

\nopagebreak[4]
\[
F_\delta'=H_0'\in\{0,1\}
\quad\text{almost everywhere on }Z.
\]

The complement of $Z$ is a countable union of disjoint open intervals. In each such
interval choose a bi-infinite sequence of subdivision points tending to both endpoints.
Every compact cell $I_j$ between consecutive subdivision points lies strictly inside
the corridor, so

\nopagebreak[4]
\[
\rho_j=\min_{I_j}\{F_\delta-\ell,u-F_\delta\}>0.
\]

Apply Lemma 3.2 on $I_j$ with mesh width smaller than
$\min\{\varepsilon,\rho_j\}$. The replacement agrees with $F_\delta$ at the cell
endpoints, remains in the corridor, and has zero--one derivative almost everywhere.
Keep $F_\delta$ on $Z$ and at all subdivision points. This defines a continuous path
$H$ on $[0,1]$.

For completeness, we verify the global regularity rather than merely the cellwise
regularity. Let $S$ be the union of $Z$ and all subdivision points. It is closed: new
accumulation points of subdivision points lie in $Z$. If $s<t$ lie in different cells,
let $r\in S$ be the right endpoint of the cell containing $s$, with the convention
$r=s$ when $s\in S$, and let $q\in S$ be the left endpoint of the cell containing $t$,
with the convention $q=t$ when $t\in S$. These conventions also remove any ambiguity
when an endpoint is adjacent to more than one cell. Since $H=F_\delta$ on $S$,

\nopagebreak[4]
\[
0\le H(t)-H(s)
\le (t-q)+F_\delta(q)-F_\delta(r)+(r-s)
\le t-s.
\]

Thus $H$ is globally nondecreasing and $1$-Lipschitz, hence absolutely continuous.
On $[0,1]\setminus S$ its derivative is zero--one by construction. On $S$, the
function $H-F_\delta$ vanishes; its derivative is therefore zero almost everywhere on
$S$. The set $S\setminus Z$ is countable, while $F_\delta'$ is zero--one almost
everywhere on $Z$. Consequently

\nopagebreak[4]
\[
H'=\mathbf 1_B
\]

almost everywhere for a measurable set $B$. The endpoint and corridor conditions show
that $B\in\mathcal B_K(a,b)$. Moreover,

\nopagebreak[4]
\[
\|H-F\|_\infty\le\varepsilon+\delta.
\]

Letting $\varepsilon,\delta\downarrow0$ gives basis paths converging uniformly to $F$.

\subsection{Weak-star convergence of the derivatives}


Suppose $H_n\to F$ uniformly and $0\le H_n',F'\le1$. For every
$\varphi\in C^1([0,1])$, integration by parts gives

\nopagebreak[4]
\[
\int_0^1\varphi(H_n'-F')
=\varphi(1)(H_n-F)(1)-\varphi(0)(H_n-F)(0)
-\int_0^1\varphi'(H_n-F),
\]

which tends to zero. Density of $C^1([0,1])$ in $L^1([0,1])$, together with the
uniform $L^\infty$ bound on the derivatives, yields $H_n'\overset{*}{\rightharpoonup}F'$.
Hence $f\lambda$ lies in the weak-star closure of the basis measures, which is
$\mathsf{bmm}(M_K(a,b))$ by \cite{berczi2026}. This completes the proof of Theorem 3.1.
\hfill\qedsymbol


\section{The full independence relaxation and all truncations}


We now allow the total mass to vary.

\begin{theorem}[interval description of all minorizing measures]
\label{thm:minorizing-interval}

For every nonempty measurable lattice-path matroid $M_K(a,b)$,

\nopagebreak[4]
\[
\mathsf{mm}(M_K(a,b))
=\{f\lambda:f\in\mathcal I_K(a,b)\}.
\]

\end{theorem}

\begin{proof}[Proof of necessity]

Let $I$ be independent. It extends to a basis $B$. For $s\le t$ in $K$,

\nopagebreak[4]
\[
\lambda(I\cap(s,t])
\le\lambda(B\cap(s,t])
=F_B(t)-F_B(s)
\le b(t)-a(s).
\]

It follows that the rank of $(s,t]$ is at most $b(t)-a(s)$. Every minorizing measure
therefore satisfies the interval inequalities. Its density also satisfies $0\le f\le1$.

\medskip
\noindent\emph{Proof of sufficiency.}


Fix $f\in\mathcal I_K(a,b)$, set $F=F_f$, and write

\nopagebreak[4]
\[
m=F(1),\qquad C(t)=t-F(t)=\int_0^t(1-f(x))\,dx.
\]

The function $C$ is nondecreasing and $1$-Lipschitz. Put $U=1-m$ and
$S=C(K)\subseteq[0,U]$. For $u\in S$, define

\nopagebreak[4]
\[
A(u)=\sup\{a(t)-F(t):t\in K,\ C(t)=u\},
\]

and

\nopagebreak[4]
\[
B(u)=\inf\{b(t)-F(t):t\in K,\ C(t)=u\}.
\]

The fibers are compact. We claim that for $u\le v$ in $S$,

\nopagebreak[4]
\[
A(u)\le B(v),\qquad A(v)-B(u)\le v-u. \tag{4.1}
\]

Take $s,t\in K$ with $C(s)=u$ and $C(t)=v$. If $u<v$, monotonicity of $C$ forces
$s<t$. The interval inequality gives

\nopagebreak[4]
\[
a(s)-F(s)\le b(t)-F(t),
\]

while monotonicity and the $1$-Lipschitz property of $a$ give

\nopagebreak[4]
\[
a(t)-b(s)\le t-s.
\]

After subtracting $F(t)-F(s)$, the latter becomes

\nopagebreak[4]
\[
a(t)-F(t)-\bigl(b(s)-F(s)\bigr)
\le C(t)-C(s)=v-u.
\]

Taking the appropriate supremum and infimum proves (4.1). If $u=v$, points in the same
fiber may occur in either order. Applying the first calculation in the forward order
and the $1$-Lipschitz calculation in the reverse order gives both inequalities with
right-hand side zero. Thus (4.1) holds for non-singleton fibers as well.

Define, for $0\le u\le U$,

\nopagebreak[4]
\[
y(u)=\sup_{v\in S}\{A(v)-\max(v-u,0)\}. \tag{4.2}
\]

Each function in the supremum is nondecreasing and $1$-Lipschitz, so the same is true
of $y$. Choosing $v=u$ gives $A(u)\le y(u)$ for $u\in S$. If $v\le u$, the first
inequality in (4.1) gives $A(v)\le B(u)$; if $v\ge u$, the second gives
$A(v)-(v-u)\le B(u)$. Hence

\nopagebreak[4]
\[
A(u)\le y(u)\le B(u)\qquad(u\in S). \tag{4.3}
\]

The endpoint fibers give

\nopagebreak[4]
\[
A(0)=B(0)=0,qquad A(U)=B(U)=R-m.
\]

For the first pair, $C(t)=0$ implies $F(t)=t$, and the interval inequality from $0$ to
$t$, together with $b(t)\le t$, forces $b(t)=t$. For the second pair, $C(t)=U$ means
that $C$ is constant from $t$ to $1$; the interval inequality and the Lipschitz boundary
conditions then squeeze both residual boundaries to $R-m$, with equality at $t=1$.
Equations (4.1)--(4.2) consequently give

\nopagebreak[4]
\[
y(0)=0,qquad y(U)=R-m. \tag{4.4}
\]

Set

\nopagebreak[4]
\[
H(t)=y(C(t)).
\]

Both $y$ and $u\mapsto u-y(u)$ are nondecreasing. Therefore $H$ and $C-H$ are
nondecreasing absolutely continuous functions. It follows that

\nopagebreak[4]
\[
0\le H'\le C'=1-f
\]

almost everywhere. For $t\in K$, (4.3) gives

\nopagebreak[4]
\[
a(t)-F(t)\le H(t)\le b(t)-F(t).
\]

Thus

\nopagebreak[4]
\[
g=f+H'
\]

satisfies $0\le g\le1$, has total mass $R$ by (4.4), and its cumulative function
$F+H$ belongs to the corridor $\mathcal P_K(a,b)$. Theorem 3.1 gives
$g\lambda\in\mathsf{bmm}(M_K(a,b))$.

Finally, $0\le f\le g$. Hence, for every measurable $X$,

\nopagebreak[4]
\[
\int_X f\,d\lambda
\le\int_X g\,d\lambda
\le r(X).
\]

Therefore $f\lambda\in\mathsf{mm}(M_K(a,b))$.
\end{proof}


\begin{corollary}[all real truncations]
\label{cor:truncations}

Let $0\le c\le R$, and let $T_cM_K(a,b)$ denote the $c$-truncation. Then

\nopagebreak[4]
\[
\mathsf{bmm}(T_cM_K(a,b))
=\left\{f\lambda:
f\in\mathcal I_K(a,b),\ \int_0^1f=c
\right\}.
\]

\end{corollary}

\begin{proof}

The rank of a truncation is $r_c(X)=\min\{r(X),c\}$ \cite{berczi2026}. Consequently,

\nopagebreak[4]
\[
\mathsf{mm}(T_cM)
=\{\alpha\in\mathsf{mm}(M):\alpha(J)\le c\}.
\]

Indeed, the forward inclusion follows from $r_c\le r$ and $r_c(J)=c$; conversely,
$\alpha(X)\le r(X)$ and $\alpha(X)\le\alpha(J)\le c$ imply
$\alpha(X)\le r_c(X)$. Taking the mass-$c$ slice and applying Theorem 4.1 gives the
claim. The endpoint cases $c=0$ and $c=R$ are included.
\end{proof}


When $c=R$, the interval inequalities with $(s,t)=(0,t)$ and $(t,1)$ recover the
two prefix inequalities, so Corollary 4.2 specializes to Theorem 3.1.

\section{Extreme and exposed points}


For $f\in\mathcal P_K(a,b)$, define

\nopagebreak[4]
\[
T_a(f)=\{t\in K:F_f(t)=a(t)\},
\]

\nopagebreak[4]
\[
T_b(f)=\{t\in K:F_f(t)=b(t)\},
\]

and $T(f)=T_a(f)\cup T_b(f)$. These sets are closed.

\begin{theorem}[contact-set criterion]
\label{thm:contact}

For $f\in\mathcal P_K(a,b)$, the following are equivalent:

\begin{enumerate}
\item $f\lambda$ is an extreme point of $\mathsf{bmm}(M_K(a,b))$;
\item $f\lambda$ is a weak-star exposed point of $\mathsf{bmm}(M_K(a,b))$;
\item $f$ belongs to $\{0,1\}$ almost everywhere on $[0,1]\setminus T(f)$.

\end{enumerate}

\end{theorem}

\begin{proof}[Proof of necessity]

By Theorem 3.1 it suffices to work in $\mathcal P_K(a,b)$. If the fractional part of
$f$ outside $T(f)$ has positive measure, then some compact interval contained in one
component of $[0,1]\setminus T(f)$ contains a positive-measure set on which
$\eta\le f\le1-\eta$ for some $\eta>0$. On the constrained times in that compact
interval, both relevant slacks have a positive uniform lower bound. Choose disjoint
equal-measure sets $E_+$ and $E_-$ in the fractional region and put
$h=\mathbf1_{E_+}-\mathbf1_{E_-}$. Its cumulative integral vanishes outside the
chosen interval. For sufficiently small $\varepsilon>0$, both
$f+\varepsilon h$ and $f-\varepsilon h$ remain in the corridor. Hence $f$ is not
extreme.

\medskip
\noindent\emph{Proof of sufficiency.}


Suppose $f=(g_1+g_2)/2$ with $g_i\in\mathcal P_K(a,b)$. Put
$d=(g_1-g_2)/2$ and $D(t)=\int_0^t d$. The box constraints force $d=0$ almost
everywhere where $f=0$ or $f=1$, hence almost everywhere off $T(f)$. At a point of
$T_a(f)$ or $T_b(f)$, equality in the averaged active constraint forces both $g_i$ to
take the same cumulative value as $f$. Therefore $D=0$ on $T(f)$. An absolutely
continuous function has derivative zero almost everywhere on its zero set, so $d=0$
almost everywhere on $T(f)$ as well. Thus $g_1=g_2=f$.

\medskip
\noindent\emph{Exposedness.}


Choose Borel probability measures $\nu_a$ and $\nu_b$ with full support on the two
nonempty compact contact sets. Define

\nopagebreak[4]
\[
q_b(x)=\nu_b(T_b(f)\cap[x,1]),\qquad
q_a(x)=\nu_a(T_a(f)\cap[x,1]).
\]

On the complement of $T(f)$, choose a zero--one representative of $f$ and set
$c_f=1$ where $f=1$ and $c_f=-1$ where $f=0$; put $c_f=0$ on $T(f)$. Then

\nopagebreak[4]
\[
w=c_f+q_b-q_a\in L^1([0,1]).
\]

For every $g\in\mathcal P_K(a,b)$, the box constraints and Fubini's theorem give

\nopagebreak[4]
\[
\int c_fg\le\int c_ff,
\]

\nopagebreak[4]
\[
\int q_bg
=\int_{T_b(f)}F_g(t)\,d\nu_b(t)
\le\int_{T_b(f)}b(t)\,d\nu_b(t)
=\int q_bf,
\]

and similarly $-\int q_ag\le-\int q_af$. If equality holds in their sum, every
nonnegative gap is zero. The first equality fixes $g=f$ off $T(f)$. Full support and
continuity force all cumulative contact equalities, and the zero-set derivative lemma
then fixes $g=f$ on $T(f)$ almost everywhere. Thus $w$ uniquely exposes $f$. If an
explicit nonzero exposing representative is desired in a degenerate case, add the
constant function $1$; the total mass is fixed on the corridor.
\end{proof}


The extreme/basic part of this theorem is analogous to classical SCLP extreme-point
criteria \cite{anderson1983,andersonnash1987}. The direct argument above is self-contained; the references provide
background rather than a claimed exact theorem-level reduction.

\begin{corollary}[integrality for a fixed corridor]
\label{cor:integrality}

Let $\dot a_K$ and $\dot b_K$ denote the intrinsic almost-everywhere derivatives on
$K$: take any Lipschitz extensions to $[0,1]$; their derivatives agree almost
everywhere on $K$. Every extreme point of $\mathsf{bmm}(M_K(a,b))$ is induced by a
basis if and only if

\nopagebreak[4]
\[
\dot a_K,\dot b_K\in\{0,1\}
\quad\text{almost everywhere on }K.
\]

\end{corollary}

\begin{proof}

Theorem 5.1 already forces zero--one values off the contact set. On the lower and upper
contact sets, the zero-set derivative lemma gives $f=\dot a_K$ and
$f=\dot b_K$, respectively, almost everywhere. This proves sufficiency.

Conversely, if $\dot b_K$ is fractional on a positive-measure subset, extend $b$ over
each complementary gap by moving first with slope $1$ and then with slope $0$, while
keeping $F=b$ on $K$. The resulting path belongs to the corridor, is zero--one off
$K$, and is in upper contact throughout $K$. Theorem 5.1 makes its derivative an
exposed point, but it is fractional on a positive-measure subset of $K$. The lower
boundary is handled analogously.
\end{proof}


\section{A bounded laminar-composition application}


This section records a useful closure consequence but is not a separate main
contribution.

Let $T$ be a finite rooted tree whose internal nodes have finite arity. Its leaves carry
pairwise disjoint atomless standard probability spaces, each identified measure
preservingly with $[0,1]$. At a leaf $\ell$, let
$K_\ell\subseteq[0,1]$ be closed and contain $0,1$, and let
$b_\ell:K_\ell\to[0,\infty)$ be an effective, nondecreasing, $1$-Lipschitz upper
capacity with $b_\ell(0)=0$. If one starts from an arbitrary capacity presentation, it
must first be replaced by its effective rank representative: for a nested family
$\mathcal D$ with raw capacity $b$, use the rank-effective capacity

\nopagebreak[4]
\[
\beta(C)=\inf_{D\in\mathcal D}\{b(D)+\mu(C\setminus D)\}
\qquad(C\in\mathcal D),
\]

and then parameterize the resulting closed chain by $K_\ell$. The contact-set statement
below is not asserted for a raw discontinuous or redundant presentation.

Let

\nopagebreak[4]
\[
C_\ell=
\left\{f_\ell:0\le f_\ell\le1,
\ \int_0^t f_\ell\le b_\ell(t)\quad(t\in K_\ell)
\right\}.
\]

Every internal node is labelled either by the direct sum of its children or by a
truncation of the matroid represented by its unique child; every truncation capacity
$c_v$ is a nonnegative real number. Consecutive direct sums may be merged and
consecutive truncations replaced by the smaller capacity. If $L(v)$ is the set of
leaves below a truncation node $v$, define

\nopagebreak[4]
\[
Q_T=\left\{(f_\ell):
f_\ell\in C_\ell,
\ \sum_{\ell\in L(v)}\int f_\ell\le c_v
\text{ for every truncation node }v
\right\}.
\]

\begin{proposition}
\label{prop:laminar-mm}

For the measurable matroid $M_T$ obtained from this tree,

\nopagebreak[4]
\[
\mathsf{mm}(M_T)=Q_T,
\]

and

\nopagebreak[4]
\[
\mathsf{bmm}(M_T)
=\left\{f\in Q_T:\sum_\ell\int f_\ell=r_T(J)\right\}.
\]

\end{proposition}

\begin{proof}

At a leaf, the effective nested rank formula says that the minorizing-measure set is
exactly $C_\ell$. Minorizing-measure sets multiply under direct sum. Under a
$c$-truncation they are intersected with the total-mass halfspace of level $c$, by the
same rank calculation used in Corollary 4.2. Induction over the finite tree gives the
claim.
\end{proof}


For the extreme-point statement, define the leaf contact set

\nopagebreak[4]
\[
T_\ell(f_\ell)
=\left\{t\in K_\ell:\int_0^t f_\ell=b_\ell(t)\right\}.
\]

We need one elementary convex lemma.

\begin{lemma}[one affine slice of a product]
\label{lem:affine-slice}

Let $C_i$ be convex sets, let $\ell_i$ be affine real functionals, and put

\nopagebreak[4]
\[
C(c)=\left\{(x_i)\in\prod_iC_i:\sum_i\ell_i(x_i)=c\right\}.
\]

A point $x=(x_i)$ is extreme in $C(c)$ if and only if every $x_i$ is extreme in its
fixed-value fiber

\nopagebreak[4]
\[
\{y\in C_i:\ell_i(y)=\ell_i(x_i)\},
\]

and at most one $x_i$ is not extreme in the whole set $C_i$.

\end{lemma}

\begin{proof}

A nontrivial symmetric perturbation inside one fixed-value fiber immediately gives a
perturbation in $C(c)$. If two coordinates admit whole-set perturbations but no
fixed-fiber perturbations, both perturbations change their affine values. Reverse one
and scale the two symmetric segments so that the changes cancel. Conversely, any
nontrivial perturbation in $C(c)$ must, after fixed-fiber extremality is imposed, change
the affine value in every perturbed coordinate. Since the changes sum to zero, at least
two coordinates are perturbed.
\end{proof}


\begin{proposition}[leafwise extreme-point recursion]
\label{prop:laminar-extreme}

For every node $v$ of the tree, a feasible density is extreme in the whole subtree
relaxation if and only if it is extreme in its fixed-total-mass fiber, if and only if
every leaf below $v$ satisfies

\nopagebreak[4]
\[
f_\ell\in\{0,1\}
\quad\text{almost everywhere on }[0,1]\setminus T_\ell(f_\ell).
\]

\end{proposition}

\begin{proof}

At a leaf, the same perturbation and zero-set derivative argument as in Theorem 5.1
shows that whole-set and fixed-mass-fiber extremality coincide and are characterized by
the displayed contact condition. At a direct-sum node, whole-set extremality is
coordinatewise, while fixed-mass extremality is governed by Lemma 6.2. The induction
hypothesis makes the lemma's extra alternative redundant. At a truncation node, an
inactive capacity leaves the original convex set locally unchanged, while an active
capacity replaces it by the corresponding fixed-mass fiber. The induction hypothesis
again identifies the two extreme-point notions.
\end{proof}


Finite laminar matroids are classically organized by direct sums, truncations, and
related basic operations \cite{fifeoxley2017}. Propositions 6.1--6.3 should therefore be viewed as a
measurable application of the preceding convex mechanism, not as a new structural
classification of laminar matroids.

\section{Scope and open boundary}


Theorems 3.1 and 4.1 give exact infinite-dimensional counterparts of familiar finite
lattice-path and positroid polytope descriptions. The finite shadow is not new: base
polytopes of lattice-path matroids have path-inequality descriptions \cite{bidkhori2012}, and general
positroid independence polytopes admit interval-type facet descriptions \cite{ohxiang2022}. The new
claim made here is limited to the atomless measurable setting, arbitrary closed $K$,
the explicit countable purification and residual-capacity extension, and all real
truncation levels. No claim of absolute global priority is made without external review.

Two natural questions remain.

\begin{enumerate}
\item Give a complete extreme-point criterion for the mass-$c$ slice of
$\mathcal I_K(a,b)$ when $c<R$. The equalities on active intervals give necessary
perturbation conditions, but nonactive slacks may converge to zero, so the obvious
finite-dimensional kernel argument is not sufficient.
\item Determine when finite laminar compositions as in Section 6 are genuinely outside
the class of measurable lattice-path matroids. The classical finite truncation of
$U_{1,2}^{\oplus3}$ is not transversal and hence not a lattice-path matroid \cite{bonindemier2006}, but
transferring this into a sharp atomless nonrepresentation theorem requires a
separate argument.

\end{enumerate}

These questions require new structure beyond the direct-sum and truncation identities
used here.

\section*{Acknowledgement of AI assistance}


OpenAI Codex was used as an AI-assisted research and writing tool for literature
discovery, exploratory checks, proof auditing, and language editing. Carptopus takes
responsibility for the mathematical claims and the released version.

\begin{thebibliography}{9}
\footnotesize
\raggedright

\bibitem{berczi2026}
K.~B\'erczi, A.~Imolay, and \'A.~Schweitzer,
\emph{Measurable Matroids: Foundations and Min--Max Theorems},
arXiv:2608.16464v1 (2026).
\href{https://arxiv.org/abs/2608.16464}{arXiv:2608.16464}.

\bibitem{bidkhori2012}
H.~Bidkhori,
\emph{Lattice Path Matroid Polytopes},
arXiv:1212.5705 (2012).
\href{https://arxiv.org/abs/1212.5705}{arXiv:1212.5705}.

\bibitem{ohxiang2022}
S.~Oh and D.~Xiang,
\emph{The facets of the matroid polytope and the independent set polytope of a positroid},
J. Combin. 13(4) (2022), 545--560.
\href{https://doi.org/10.4310/JOC.2022.V13.N4.A5}{doi:10.4310/JOC.2022.V13.N4.A5};
\href{https://arxiv.org/abs/1701.08483}{arXiv:1701.08483v3}.

\bibitem{anderson1983}
E.~J. Anderson, P.~Nash, and A.~F. Perold,
\emph{Some Properties of a Class of Continuous Linear Programs},
SIAM J. Control Optim. 21 (1983), 758--765.
\href{https://doi.org/10.1137/0321046}{doi:10.1137/0321046}.

\bibitem{andersonnash1987}
E.~J. Anderson and P.~Nash,
\emph{Linear Programming in Infinite-Dimensional Spaces: Theory and Applications},
Wiley, 1987, Section 7.4.

\bibitem{amarcellina1994}
M.~Amar and A.~Cellina,
\emph{On passing to the limit for non-convex variational problems},
Asymptotic Anal. 9 (1994), 135--148.
\href{https://doi.org/10.3233/ASY-1994-9203}{doi:10.3233/ASY-1994-9203}.

\bibitem{carlota2016}
C.~Carlota, S.~Ch\'a, and A.~Ornelas,
\emph{A pointwise constrained version of the Liapunov convexity theorem for vectorial linear first-order control systems},
J. Differential Equations 261 (2016), 296--318.
\href{https://doi.org/10.1016/j.jde.2016.03.009}{doi:10.1016/j.jde.2016.03.009}.

\bibitem{fifeoxley2017}
T.~Fife and J.~Oxley,
\emph{Laminar Matroids},
European J. Combin. 62 (2017), 206--216.
\href{https://doi.org/10.1016/j.ejc.2017.01.002}{doi:10.1016/j.ejc.2017.01.002}.

\bibitem{bonindemier2006}
J.~E. Bonin and A.~de Mier,
\emph{Lattice Path Matroids: Structural Properties},
European J. Combin. 27 (2006), 701--738.
\href{https://doi.org/10.1016/j.ejc.2005.01.008}{doi:10.1016/j.ejc.2005.01.008}.

\end{thebibliography}

\end{document}
